%!TEX root = ../../main.tex
\section{공간 경계 $D$}
\label{sec:eng-spatial}

시간적으로 가까운 킬이 다른 전투에 속할 수 있다. 두 라인에서 동시에 벌어지는 전투가 그 예이다. 따라서 에피소드는 어떤 두 킬도 $D$보다 멀지 않도록 나뉜다. 시간 순으로 각 킬은 구성원 모두가 자신으로부터 $D$ 안에 있는 첫 그룹에 합류하거나 새 그룹을 만든다.

$D$는 시간 규칙이 쓰지 않는 신호에서 추정한다. 연속한 두 킬이 같은 챔피언을 킬러·희생자·어시스트로 공유하는지이다. 한 챔피언은 동시에 두 곳에 있을 수 없으므로, $\hat G$ 이내의 쌍 가운데 거리의 함수로 본 챔피언 공유율은 ``같은 전투''가 어디서 끝나는지를 그린다. $D$는 그 비율이 0.5로 떨어지는 거리, 곧 두 오류의 비용이 같을 때 가장 적은 쌍을 잘못 분류하는 결정점이다. $\hat G$ 이내의 \numPairsDef{} 연속 쌍에서 $\hat D = \constDest$ 단위(경기 부트스트랩 95\% 구간 \constDciLo{}--\constDciHi{})이다. 규칙이 연결하는 쌍은 \shareLinkedPairs{}\%에서 챔피언을 공유하고, $\hat G$ 이내이지만 규칙이 분리하는 쌍은 \shareSeparatedPairs{}\%에서만 공유한다.

공유율은 거리 구간별로 계산되고 $\hat D$는 0.5를 사이에 둔 두 구간 중심(3{,}750과 4{,}500 단위) 사이의 선형 보간이므로, 18 단위 폭의 부트스트랩 구간은 두 구간 평균의 표본 변동만을 재고 750 단위 간격에서 비율이 선형이라는 가정의 오차는 보지 못한다. 따라서 $\hat D$의 해상도는 이 구간 간격 수준이다. 라인 안에서 경계는 비등방적이어서 라인 방향의 교차점이 라인 가로 방향의 1.5--1.6배이지만, 타원 경계는 파일럿에서 0.16\%의 쌍만을 다시 분류했고 지형 경로 거리는 직선 거리보다 나은 분리를 주지 않았으므로 등방 규칙을 유지한다. 세 패치의 $D$는 \constDpatchA{}, \constDpatchB{}, \constDpatchC{} 단위로 합동 값의 10\% 안에 있다. 검출기는 $D = \constD$ 단위로 실행되었다.
